% Fluxo visual: matriz A → matriz dos cofatores → adjunta (transposta dos cofatores) → inversa.
% Caso 2x2: A = [[1,2],[3,4]], det A = -2.
% Usado em: capitulo_02_propriedades-e-aplicacoes.md (seção 5)
\begin{tikzpicture}[
  every node/.style={font=\small},
  caixa/.style={
    rectangle,
    rounded corners=4pt,
    draw=eleveBlue,
    fill=eleveBlue!8,
    minimum width=28mm,
    minimum height=22mm,
    inner sep=4pt,
    align=center,
    font=\small
  },
  flecha/.style={->,thick,>=stealth,eleveAccent},
  rotulo/.style={font=\scriptsize,eleveAccent}
]
  % --- A ---
  \node[caixa] (A) at (0,0) {$A$\\[1pt] linha 1: $(1,2)$\\ linha 2: $(3,4)$\\[2pt] $\det A = -2$};

  % --- cof(A) ---
  \node[caixa] (cof) at (4.5,0) {$\mathrm{cof}(A)$\\[1pt] linha 1: $(4,-3)$\\ linha 2: $(-2,1)$};

  % --- adj(A) ---
  \node[caixa] (adj) at (9,0) {$\mathrm{adj}(A) = (\mathrm{cof}\,A)^T$\\[1pt] linha 1: $(4,-2)$\\ linha 2: $(-3,1)$};

  % --- inversa ---
  \node[caixa,fill=eleveAccent!10,draw=eleveAccent] (inv) at (13.5,0) {$A^{-1} = \dfrac{1}{\det A} \cdot \mathrm{adj}(A)$\\[1pt] linha 1: $(-2,1)$\\ linha 2: $(1{,}5;\,-0{,}5)$};

  % Setas com rótulos
  \draw[flecha] (A) -- (cof) node[rotulo,midway,above,align=center] {calcular\\ cofatores};
  \draw[flecha] (cof) -- (adj) node[rotulo,midway,above,align=center] {transpor};
  \draw[flecha] (adj) -- (inv) node[rotulo,midway,above,align=center] {dividir por\\ $\det A$};
\end{tikzpicture}
